\section{Conclusions}
\label{sec:conclusions}

\subsection{Résumé des contributions}
\label{subsec:contributions}

Ce rapport a abordé un problème documenté et économiquement significatif dans l'industrie du bâtiment :
la cartographie manuelle des points BACnet coûte 0,5 à 1,5\% de la valeur de construction par bâtiment,
bloque l'interopérabilité logicielle, et est identifiée dans la littérature évaluée par les pairs et les
rapports de politique du DOE comme un obstacle majeur à l'innovation dans l'automatisation des bâtiments.
Le même problème empêche directement l'inscription à grande échelle des bâtiments commerciaux dans les
programmes de Centrale électrique virtuelle et d'Alternative non filaire, car chaque bâtiment nécessite
un effort d'intégration personnalisé, par projet, avant de pouvoir participer aux services réseau.
Aucune solution automatisée largement déployée pour ce problème n'existait avant ce travail.

SI-Mapper démontre que l'IA agentique peut automatiser la création de modèles sémantiques de bâtiment
ASHRAE 223P à partir des entrées brutes disponibles sur tout projet de bâtiment réel : plans d'ingénierie
HVAC et exportations de points BACnet.
Le système combine la reconnaissance multimodale de plans, utilisant des modèles de fondation capables
d'interpréter directement les symboles schématiques HVAC à partir d'images, avec la cartographie
automatisée des points BACnet vers les équipements et la génération d'ontologie ASHRAE 223P formelle dans
un pipeline unique de bout en bout.
L'agent valide sa propre sortie d'ontologie par autocorrection itérative, en bouclant jusqu'à ce que le
fichier Turtle (TTL) généré s'exécute sans erreurs et que le graphe résultant se charge proprement dans
le magasin sémantique Neo4j.
Ce pipeline remplace un flux de travail qui nécessitait auparavant des jours de travail d'ingénierie HVAC
expert.
Deux expériences sur des installations réelles d'unités de traitement d'air confirment la viabilité de
bout en bout du pipeline : le Système 1 (21 composants, 44 points BACnet, 295 nœuds et 949 relations dans
le graphe résultant) a été complété en 17 minutes 11 secondes à un coût estimé de \$11.30 USD ; le
Système 3 (27 composants, 34 points BACnet, boucle de récupération de chaleur plus complexe) a été
complété en 12 minutes 24 secondes à \$10.85 USD.
Dans les deux exécutions, l'agent a identifié et corrigé de manière autonome les erreurs de code
d'ontologie, trois itérations pour le Système 1, une pour le Système 3, sans intervention humaine requise.
Les deux sessions ont utilisé Gemini 3.1 Pro Preview et sont restées bien dans la fenêtre de contexte
du modèle tout au long.

Le processus de développement lui-même est une contribution.
Neuf itérations expérimentales, des simples appels de fonctions LangChain en 2025 jusqu'à l'architecture
finale basée sur les compétences sur Gemini 3.1 et Claude Sonnet 4.6 en 2026, constituent un registre
systématique de la façon dont les architectures d'IA agentique pour les tâches techniques spécifiques
à un domaine ont évolué avec l'amélioration des capacités des modèles de fondation.
La progression révèle une leçon architecturale contre-intuitive : les architectures plus simples sont
devenues plus viables, et non moins, à mesure que les modèles s'amélioraient.
Les conceptions multi-agents avec des sous-agents dédiés Planificateur, Acteur et Réviseur, initialement
nécessaires pour maintenir le pipeline sur la bonne voie, sont devenues inutiles une fois que les modèles
ont atteint une capacité suffisante de suivi d'instructions.
L'agent maître unique final basé sur les compétences a surpassé chaque architecture multi-agents testée
tout en étant considérablement plus simple à maintenir, déboguer et étendre.

La pile technologique assemblée pour SI-Mapper, Google ADK, CopilotKit, Neo4j avec Neosemantics
(n10s), et Sigma.js, fournit une architecture de référence réutilisable pour les applications d'IA
agentique nécessitant une observabilité en temps réel par l'opérateur, un stockage de graphes sémantiques,
et une visualisation interactive de graphes.
La combinaison d'ADK pour l'orchestration d'agents, de CopilotKit pour le protocole de streaming
Agent-to-UI, et de Neo4j pour le stockage natif de graphes RDF n'est pas une combinaison évidente ;
les décisions de conception qui y ont conduit sont documentées dans la section~\ref{sec:implementation}
comme guide pratique pour les futurs implémenteurs.

\subsection{Limitations}
\label{subsec:limitations}

\textbf{Dépendance aux modèles.}
La performance du système est directement couplée aux capacités et à la disponibilité de modèles de
fondation spécifiques, Gemini 3.1 et Claude Sonnet 4.6 au moment de la rédaction.
Les changements d'API de modèle, les changements de tarification ou les régressions de capacité dans
les futures versions de modèle pourraient affecter la fiabilité du pipeline sans aucun changement dans
le code source de SI-Mapper.
L'architecture des compétences atténue quelque peu ce risque en encapsulant les connaissances de domaine
dans des ensembles d'instructions en texte brut qui peuvent être mis à jour indépendamment des versions
de modèle, mais la dépendance sous-jacente aux fournisseurs de modèles externes demeure.

\textbf{Portée de validation limitée.}
La portée expérimentale actuelle couvre deux systèmes HVAC réels (Système 1 et Système 3), deux unités
de traitement d'air tirées du même contexte de projet et suivant les mêmes conventions de dénomination
BACnet en français québécois.
Les deux expériences ont réussi, l'agent a produit un fichier Turtle ASHRAE 223P valide en une seule
session dans chaque cas, mais deux systèmes d'un seul entrepreneur et d'une seule région sont insuffisants
pour établir une confiance statistique dans la précision ou la généralisabilité.
Aucune exécution échouée n'a été observée, bien que des erreurs mineures de code d'ontologie (imports
incorrects, paramètres de constructeur manquants, syntaxe d'opérateur) se soient produites et aient été
corrigées de manière autonome dans tous les cas.
Si ces erreurs seraient récupérables sur des systèmes avec des styles de plans significativement
différents, des listes de points plus denses ou des topologies d'équipements plus inhabituelles n'est
pas encore connu.

\textbf{Spécificité de domaine.}
L'ensemble actuel de compétences et d'outils est adapté aux unités de traitement d'air HVAC avec des
points BACnet dans des conventions de dénomination en langue française, reflétant le contexte
d'installation québécois où ce travail a été développé.
Étendre SI-Mapper à d'autres types de systèmes de bâtiment (boîtes VAV, installations de refroidisseurs,
installations de chaudières), d'autres protocoles de contrôle (Modbus, KNX, BACnet/IP vs.\ BACnet MSTP),
ou des conventions de dénomination de points BACnet d'autres régions et entrepreneurs nécessiterait un
développement et une validation supplémentaires des compétences.
L'architecture supporte cette extension naturellement, les nouvelles compétences sont des ensembles
d'instructions en texte brut, mais les connaissances de domaine nécessaires pour écrire des compétences
précises sont non triviales.

\textbf{Maturité d'ASHRAE 223P.}
La norme est actuellement en deuxième examen public consultatif et n'est pas encore ratifiée.
La logique de génération d'ontologie dans SI-Mapper est calibrée par rapport au projet consultatif actuel.
Des modifications de la norme lors de la ratification, en particulier pour les types de points de
connexion, les définitions de médiums ou la syntaxe de liaison de télémétrie, pourraient nécessiter des
mises à jour des compétences de génération d'ontologie et de la bibliothèque Python ASHRAE 223P sous-jacente
utilisée pour produire la sortie TTL.

\subsection{Travaux futurs}
\label{subsec:future-work}

\textbf{Validation en production à grande échelle.}
L'étape suivante la plus importante est de continuer à développer SI-Mapper et de tester le modèle sur
un corpus plus large de bâtiments réels, avec une cible de 20 à 50 systèmes HVAC couvrant plusieurs types
de systèmes, entrepreneurs et conventions de dénomination.
Cela établirait une confiance statistique dans la précision et les économies de coûts, identifierait les
cas limites non rencontrés dans l'ensemble de développement actuel, et produirait les données de
référence nécessaires pour une publication évaluée par les pairs.

\textbf{Support multi-systèmes.}
Les compétences actuelles couvrent les unités de traitement d'air.
Un système viable en production aurait besoin de compétences pour les boîtes terminales VAV, les
installations de refroidisseurs, les installations de chaudières, les systèmes de pompe à chaleur et
les unités de ventilation mécanique.
Chaque type de système a ses propres conventions de plans, ses schémas de dénomination de points et ses
règles de topologie ASHRAE 223P.
L'architecture des compétences rend cette extension gérable : chaque type de système nécessite un nouveau
fichier de compétence plutôt qu'un nouvel agent ou une refonte du pipeline.

\textbf{Métriques de qualité automatisées.}
L'évaluation actuelle des modèles ASHRAE 223P générés est qualitative.
Un cadre d'évaluation en production inclurait des outils de comparaison automatisés : différenciation de
graphes par rapport à un modèle de référence créé par un humain, scores quantitatifs pour l'exhaustivité
(fraction des nœuds d'équipements attendus présents), la précision (fraction des références BACnet
correctement assignées), et la correction de topologie (fraction des connexions dirigées correspondant
au plan source).
Ces métriques rendraient possible le suivi des régressions à travers les mises à jour de modèles et
les révisions de compétences sans inspection manuelle.

\textbf{Support multi-langues.}
Adapter les compétences de cartographie des points BACnet pour reconnaître les conventions de dénomination
anglaises, espagnoles et autres au-delà du français élargirait significativement la géographie applicable
de SI-Mapper.
Les compétences de cartographie sont le lieu principal où les connaissances sur les conventions de
dénomination sont encodées ; l'ajout de variantes de langues est une tâche d'écriture de compétences
plutôt qu'un changement architectural.

\textbf{Projet pilote de démonstration en bâtiment réel.}
L'étape de validation à court terme ayant le plus grand impact est de connecter un modèle ASHRAE 223P
généré à des points BACnet en direct sur un bâtiment réel et d'opérer sur des données réelles.
Le projet pilote se déroulerait en deux étapes : premièrement, lier les nœuds d'ontologie à leurs
homologues physiques et vérifier que le graphe sémantique peut être utilisé pour interroger des lectures
de capteurs en direct, confirmant que l'interface d'ontologie ASHRAE 223P est suffisante pour la
surveillance des bâtiments ; deuxièmement, émettre de simples commandes de contrôle superviseur via
la même interface.
Cela validerait l'hypothèse centrale de SI-Mapper : que le modèle sémantique généré est une interface
opérationnelle utilisable pour surveiller et contrôler les bâtiments, et non simplement un artefact de
documentation amélioré.
